\chapter{Происхождение гауссова опорного состояния из родительского функционала}
\label{ch:v9-gaussian-reference-state-parent-origin}

\section{Минимальная stationary architecture}

Common covariance carrier предыдущего гейта требует опорное состояние
$\mathcal N(0,I_{10})$. Минимальная continuous Markov realization задаётся
изотропным Ornstein--Uhlenbeck процессом
\begin{equation}
 dX_t=-B X_t\,dt+\sqrt{2D}\,dW_t,
 \qquad B=\gamma I_{10},\quad D=\delta I_{10},
 \qquad \gamma,\delta>0.
 \label{eq:v9-gaussian-reference-ou-process}
\end{equation}
Его stationary covariance $S$ удовлетворяет Lyapunov equation
\begin{equation}
 BS+SB^T=2D.
 \label{eq:v9-gaussian-reference-lyapunov}
\end{equation}
Для устойчивого drift это уравнение имеет единственное решение. Связь
stationary covariance с уравнением Ляпунова стандартна для многомерной
линейной динамики Ланжевена \cite{v9-gaussian-reference-godreche2018}.

Пространство symmetric covariance directions имеет dimension
\begin{equation}
 \dim\operatorname{Sym}_{10}=\frac{10\cdot11}{2}=55.
 \label{eq:v9-gaussian-reference-symmetric-dimension}
\end{equation}
При $B=I_{10}$ induced Lyapunov map на нём равен $2I_{55}$, поэтому
\begin{equation}
 \rank\mathcal L_B=55,
 \qquad \operatorname{nullity}\mathcal L_B=0.
 \label{eq:v9-gaussian-reference-lyapunov-rank}
\end{equation}
Итак, при уже фиксированных $B,D$ stationary covariance действительно
единственна.

\section{Свободное отношение diffusion/drift}

Точное решение изотропной задачи равно
\begin{equation}
 S_{\gamma,\delta}=\frac{\delta}{\gamma}I_{10}.
 \label{eq:v9-gaussian-reference-covariance-family}
\end{equation}
Каждый член семейства stable, isotropic и reversible, поскольку
$BD=DB^T$. В частности, существуют два exact witnesses
\begin{equation}
 (\gamma,\delta)=(1,1)\Longrightarrow S=I_{10},
 \qquad
 (\gamma,\delta)=(1,2)\Longrightarrow S=2I_{10}.
 \label{eq:v9-gaussian-reference-two-witnesses}
\end{equation}
Оба проходят одинаковые structural tests, но
\begin{equation}
 \Tr(2I_{10}-I_{10})=10.
 \label{eq:v9-gaussian-reference-witness-separation}
\end{equation}
Следовательно, unit covariance требует дополнительного relation
$\delta=\gamma$; оно не следует из stationarity, isotropy, reversibility
или uniqueness.

Более того,
\begin{equation}
 (\gamma,\delta)\longmapsto(s\gamma,s\delta)
 \qquad\Longrightarrow\qquad
 \frac{s\delta}{s\gamma}=\frac{\delta}{\gamma}.
 \label{eq:v9-gaussian-reference-common-rescaling}
\end{equation}
В logarithmic coefficients covariance-map имеет tangent
\begin{equation}
 J_{\rm cov}=\begin{pmatrix}-1&1\end{pmatrix},
 \qquad \rank J_{\rm cov}=1,
 \qquad \operatorname{nullity}J_{\rm cov}=1.
 \label{eq:v9-gaussian-reference-ratio-orbit}
\end{equation}

\section{Physical-origin verdict}

Endpoint QMS фиксирует finite-state relaxation architecture, но не даёт
canonical additive Gaussian diffusion amplitude $\delta$ на новом
десятимерном real carrier. Выбор $D=B$ поэтому является conditional
fluctuation--dissipation input, а не следствием старого родительского функционала.

\begin{theorem}[Conditional Gaussian опорное состояние]
Для любых fixed $\gamma,\delta>0$ OU architecture имеет unique centered
Gaussian стационарное состояние с covariance $(\delta/\gamma)I_{10}$. Unit
опорное состояние существует при $\delta=\gamma$.
\end{theorem}

\begin{nogocriterion}[Провал parent-origin unit covariance]
Поскольку два structural-equivalent witnesses дают $I_{10}$ и $2I_{10}$,
unit reference covariance не выводится из OU architecture. Relation
$\delta=\gamma$ нельзя вводить вручную как ``естественную нормировку'':
она и есть искомый reference-scale datum.
\end{nogocriterion}

\begin{protocol}\begin{enumerate}
\item Gaussian stationary architecture: \textbf{$1/1$};
\item covariance-space dimension: \textbf{$55$};
\item Lyapunov rank/nullity: \textbf{$55/0$};
\item conditional unit опорное состояние: \textbf{$1/1$};
\item physical reference-state origin: \textbf{$0/1$};
\item physical-origin packages: \textbf{$0/2$};
\item physical reopening: \textbf{нет}.
\end{enumerate}\end{protocol}

Машинный сертификат сохранён в
\path{s2t_v9_physical_reopening_gaussian_reference_state_parent_origin_gate_results.json}.

\begin{thebibliography}{9}
\bibitem{v9-gaussian-reference-godreche2018}
C.~Godr\`eche, J.-M.~Luck,
\emph{Characterising the nonequilibrium stationary states of
Ornstein--Uhlenbeck processes}, Journal of Physics A 52 (2019), 035002;
arXiv:1807.00694.
\end{thebibliography}